\documentclass[12pt]{article}
\usepackage[utf8]{inputenc}
\usepackage{listings}
\usepackage{xcolor}
\usepackage{geometry}
\usepackage{parskip}

\geometry{margin=2.5cm}

\definecolor{codebg}{RGB}{245,245,245}
\definecolor{keyword}{RGB}{0,112,0}
\definecolor{comment}{RGB}{128,128,128}
\definecolor{string}{RGB}{180,0,0}

\lstdefinestyle{pythonstyle}{
    backgroundcolor=\color{codebg},
    basicstyle=\ttfamily\small,
    keywordstyle=\color{keyword}\bfseries,
    stringstyle=\color{string},
    commentstyle=\color{comment}\itshape,
    language=Python,
    frame=single,
    rulecolor=\color{gray!40},
    breaklines=true,
    showstringspaces=false,
    tabsize=4,
    captionpos=b,
}

\title{\textbf{Wind Turbine Blade Fault Classification\\Using a Simple CNN}\\[0.5em]
\large Step 4 --- Show Sample Images}
\author{Amreen Batool}
\date{}

\begin{document}
\maketitle

\section*{Step 4 --- Show Sample Images}

\begin{lstlisting}[style=pythonstyle]
basic_datagen = ImageDataGenerator(rescale=1./255)

sample_data = basic_datagen.flow_from_directory(
    train_dir,
    target_size=(128, 128),
    batch_size=9,
    class_mode='binary',
    shuffle=True
)

images, labels = next(sample_data)

index_to_class = {v: k for k, v in sample_data.class_indices.items()}

plt.figure(figsize=(10, 8))
for i in range(9):
    plt.subplot(3, 3, i + 1)
    plt.imshow(images[i])
    plt.title(index_to_class[int(labels[i])])
    plt.axis("off")
plt.tight_layout()
plt.show()
\end{lstlisting}

\subsection*{Explanation}

\begin{itemize}
    \item It loads 9 images from the training set.
    \item It displays them in a $3 \times 3$ grid.
    \item The title above each image shows whether it is \texttt{Faulty} or \texttt{Healthy}.
\end{itemize}

\subsection*{Figure to Include}

\textbf{Figure 1: Sample Images from the Dataset}

\noindent This figure visually shows examples of both healthy and faulty turbine blade images before any model training begins.

\end{document}
